Large language models transform hidden states through many layers, yet the dynamical mechanism organizing this internal evolution remains obscure. Here we reveal an experimentally grounded picture of transformer inference as competition between trajectories in a high-dimensional layer-state space. In representative Qwen, Llama and Gemma checkpoints, row-permuted vocabulary projections preserved hidden-state geometry almost as well as the learned projection, and matched Gaussian projections often preserved more. This rejects learned vocabulary-row identity as the source of the observed organization and identifies the ordered hidden-state trajectory as the primary mechanism object. Within checkpoint-specific functional windows, a low-rank coordinate, \(\Delta U\), captured the changing relative support for leading and competing candidate trajectories. Real layer order was more directionally aligned and less curved or detoured than order-shuffled controls. Substantial residual motion excluded exact-geodesic dynamics, yet carried information complementary to the reference-aligned component. The resulting picture is a geodesic-biased competitive flow: inference advances along an organized global direction while structured off-direction motion contributes to candidate-specific processing. Trajectory-informed actions produced target-specific internal changes beyond fixed and matched shuffled policies, linking the observed geometry to an operational consequence. This dynamical organization recurred across all three checkpoints, while its magnitude and the relative strength of particular coordinates varied by checkpoint. The experiments therefore reveal a reproducible but not yet complete mechanism-level picture of internal transformer inference. They do not yet identify a full dynamical law or its behavioural consequences.

